\title{Cobweb Forests Are Universal Approximators}
\author{Draft --- working notes}
\date{\today}

\begin{document}
\maketitle

\begin{abstract}
Cobweb learns a hierarchy of probabilistic concepts one example at a time, without
gradients, by choosing among four operators with category utility.  We ask whether
Cobweb can represent and learn \emph{any} continuous function, and answer the question
for single trees and for \emph{Cobweb forests}, networks in which the predictions of some
trees are the inputs of others.  (1)~One Cobweb tree is a universal approximator.
(2)~Forests in which every tree reads at most two inputs are universal approximators, and
one input per tree is not enough.  (3)~Forests exploit composition: if the target is a
composition of $k$-input maps, a forest whose trees have $G$ leaves each approximates it
with error of order $G^{-1/k}$, an exponent independent of the number $n$ of inputs,
while the bound for one tree is $G^{-1/n}$.  (4)~Cobweb never deletes what it stores, so the tree it learns is a
consistent regressor whatever hierarchy category utility builds, and a learned forest
converges to the forest of conditional expectations its trees are asked to predict; we
characterise which of these assignments recover the target.  Every theorem predicts a
measurable quantity---a rate, a floor, an identity---and we measure each one on real
Cobweb trees; the measurements agree with the predictions, the one exception being a
single tree in eight dimensions, which is still pre-asymptotic at the largest grid we can
build.  We close with a recipe for
building a forest from data.
\end{abstract}

\tableofcontents

%=====================================================================
\section{Introduction}\label{sec:intro}
%=====================================================================

A model class is a \emph{universal approximator} if, for every continuous function on a
compact set and every tolerance, some member of the class is uniformly within that
tolerance.  Neural networks with one hidden layer are universal \citep{cybenko1989,hornik1989},
and the Kolmogorov--Arnold theorem \citep{kolmogorov1957,arnold1957}, which writes every
continuous function of $n$ variables through one-variable functions and addition, has
recently inspired an architecture with learnable one-variable units
\citep{liu2024kan}.  Universality says what a model \emph{can} represent; it is the first
thing one wants to know about a new model class.

Cobweb \citep{fisher1987} is a model class of a different kind.  It sorts each example
down a tree of probabilistic concepts and, at every node, takes whichever of four
operators (add to the best child, create a child, merge two children, split a child) most
increases \emph{category utility}, the information the concepts carry about the features
\citep{corter1992,maclellan2026rational}.  It learns incrementally, in closed form, and
without gradients.  Cobweb trees have been wired together informally before, with one
tree reading what another has learned \citep{maclellan2021conv,singaravadivelan2026trellis};
we call such a network a \emph{Cobweb forest} and give it a theory.  This paper asks:
\begin{enumerate}[label=\textbf{Q\arabic*.},leftmargin=*,noitemsep]
\item Is a Cobweb tree a universal approximator?  Is a Cobweb forest, and with how many
  inputs per tree?
\item What does composing trees buy?
\item Does Cobweb's own learning procedure learn the function, for a tree and for a forest?
\end{enumerate}

\begin{table}[t]
\centering\small
\begin{tabular}{@{}p{2.6cm}p{4.0cm}p{4.0cm}p{4.4cm}@{}}\toprule
 & multilayer perceptron & Kolmogorov--Arnold network & Cobweb forest\\\midrule
theorem & universal approximation \citep{cybenko1989,hornik1989} & Kolmogorov--Arnold \citep{kolmogorov1957} & \cref{thm:univ-tree,thm:univ-forest,thm:compose}\\
unit & neuron: fixed activation of a weighted sum & learnable spline on an edge & Cobweb tree: a concept hierarchy over its inputs\\
composition & layers of sums & sums of one-variable functions & trees reading trees, at most two inputs each\\
learning & gradient descent & gradient descent & Cobweb's own operators; no gradient (\cref{thm:learn-tree,thm:learn-forest,thm:learn-rate})\\
$k$-compositions & dimension-free rates for deep nets \citep{poggio2017} & dimension-free rate, smooth case \citep{liu2024kan} & $G^{-1/k}$ (\cref{thm:compose}), learned: $N^{-2/(k+4)}$ (\cref{thm:learn-rate})\\
\bottomrule
\end{tabular}
\caption{Three universal approximators and the theorems behind them.  This paper supplies
the right-hand column.}
\label{tab:three}
\end{table}

\paragraph{Answers.}
\begin{itemize}[leftmargin=*,noitemsep]
\item One tree is universal (\cref{thm:univ-tree}): a tree whose leaves hold the target at
  the points of a fine net, predicting with a small variance floor, is uniformly accurate.
\item Forests of two-input trees are universal (\cref{thm:univ-forest}); forests of
  one-input trees are not (\cref{prop:fanin-main}).  The proof composes trees that add and
  multiply, and uses the Stone--Weierstrass theorem; the Kolmogorov--Arnold theorem gives a
  second construction with exactly $4n^2+4n$ trees.
\item Composition beats dimension (\cref{thm:compose}): for targets composed of $k$-input
  maps a forest attains error $O(G^{-1/k})$ with $G$ leaves per tree, against the single-tree
  bound $O(G^{-1/n})$.  This is the sense in which tree--tree interactions are worth having.
\item Cobweb's learning is consistent (\cref{thm:learn-tree}), and a learned forest
  converges to an explicit limit (\cref{thm:learn-forest}), at a rate set by its fan-in when
  each tree's outcome is sufficient (\cref{thm:learn-rate}).  Whether the limit is the target
  depends on what each tree is asked to predict (\cref{prop:tasks}), and no tree can tell from
  its own data what it should predict (\cref{prop:local}).
\end{itemize}
Each answer comes with a measurement on real Cobweb trees (\cref{sec:measure},
\cref{tab:equiv}).  \Cref{sec:recipe} turns the theory into a recipe for building a forest,
and \cref{sec:cases} executes the proofs on concrete cases: a forest built literally as in
the proof of universality, and forests compiled from physics formulas.

%=====================================================================
\section{Cobweb trees and forests}\label{sec:forests}
%=====================================================================

\paragraph{A tree as a function.}
A Cobweb tree stores examples $(x,y)$, $x\in\R^d$.  Every concept $c$ keeps the number
$n_c$ of examples below it and their mean $(\mu_{c,x},\mu_{c,y})$ and variance; with exact
matching, a leaf holds $n_c$ copies of one example, so its spread is zero and its Gaussian
has the variance floor $h^2$ alone.  To predict $y$ at $x$, the tree mixes the means of its
leaves, weighting each by its count and its likelihood of $x$:
\begin{equation}\label{eq:tree}
  \hy(x)=\frac{\sum_{c\text{ leaf}}n_c\,e^{-\|x-\mu_{c,x}\|^2/2h^2}\,\mu_{c,y}}
              {\sum_{c\text{ leaf}}n_c\,e^{-\|x-\mu_{c,x}\|^2/2h^2}}.
\end{equation}
(The Gaussian normalising constants are equal for all leaves and cancel.)  This
\emph{count-weighted leaf mixture} is one of the prediction rules of the implementation;
\cref{app:model} defines all of them and says which results cover which rule.

\paragraph{Forests.}
A \emph{Cobweb forest} is a directed acyclic graph of Cobweb trees with one output tree.
Each tree reads some raw coordinates of $x$ and the \emph{messages} of its parents: their
\emph{activations}, the predictions \eqref{eq:tree} of their outcomes, or their
\emph{labels}, the concepts the input sorts to.  The forest's output is the output tree's
activation.  The \emph{fan-in} of a tree is the number of its inputs.  Every tree learns
with Cobweb's own procedure on examples made of its inputs and its \emph{outcome}, the
variable it is asked to predict (\cref{def:forest,def:task}).

%=====================================================================
\section{Representation}\label{sec:repr-main}
%=====================================================================

\subsection{One tree}

\begin{theorem}[One Cobweb tree is a universal approximator]\label{thm:univ-tree}
Let $K\subset\R^d$ be compact, $g:K\to\R$ continuous and $\varepsilon>0$.  There are a
finite set $S\subseteq K$ and $h>0$ such that the depth-one Cobweb tree whose leaves hold
the examples $(s,g(s))$, $s\in S$, satisfies $|\hy(x)-g(x)|<\varepsilon$ for every $x\in K$.
\end{theorem}

\begin{proof}
Let $M=\max_K|g|$ (finite, since $K$ is compact).  Since $g$ is uniformly continuous on
$K$, there is $\rho>0$ with $|g(x)-g(z)|<\varepsilon/2$ whenever $x,z\in K$ and
$\|x-z\|<\rho$.  Since $K$ is compact, it has a finite $\rho/2$-net: a finite $S\subseteq K$
such that every $x\in K$ is within $\rho/2$ of some $s\in S$.  Give every leaf count one.
By \eqref{eq:tree}, $\hy(x)=\sum_{s\in S}w_s(x)g(s)$ with
$w_s(x)=a_s(x)/\sum_{t\in S}a_t(x)$ and $a_s(x)=e^{-\|x-s\|^2/2h^2}$.

Fix $x\in K$ and $s_0\in S$ with $\|x-s_0\|<\rho/2$, so $\sum_ta_t(x)\ge a_{s_0}(x)>
e^{-\rho^2/8h^2}$.  Call $s$ \emph{far} if $\|x-s\|\ge\rho$; then $a_s(x)\le e^{-\rho^2/2h^2}$,
so the far leaves have total weight
\[
  \sum_{s\text{ far}}w_s(x)\;\le\;|S|\,e^{-\rho^2/2h^2}\,e^{\rho^2/8h^2}=|S|\,e^{-3\rho^2/8h^2}.
\]
Since $\sum_sw_s(x)=1$,
\[
  |\hy(x)-g(x)|\le\sum_{s\text{ near}}w_s(x)\,|g(s)-g(x)|+\sum_{s\text{ far}}w_s(x)\,|g(s)-g(x)|
  <\frac\varepsilon2+2M|S|\,e^{-3\rho^2/8h^2},
\]
where the near leaves contribute less than $\varepsilon/2$ because each is within $\rho$ of
$x$.  Choose $h>0$ so small that $2M|S|e^{-3\rho^2/8h^2}\le\varepsilon/2$; the bound is then
below $\varepsilon$ for every $x\in K$.
\end{proof}

The proof uses only that $K$ is compact and $g$ continuous; it works verbatim for any
positive counts, with $|S|$ replaced by $|S|\max_sn_s/\min_sn_s$.  For a Lipschitz $g$ a
regular grid gives an explicit rate:

\begin{proposition}[Rate; \cref{thm:single}]\label{prop:rate}
Let $g$ be $L$-Lipschitz on $K$, which lies in a cube $Q$ of side $w$, and let $\tilde g$ be
its McShane extension to $Q$ (also $L$-Lipschitz).  The depth-one tree whose $G=m^d$ leaves
hold $(\gamma,\tilde g(\gamma))$ at the cell centres $\gamma$ of an $m^d$ grid on $Q$, with
variance floor $h=\beta w/m$, predicts with the count-weighted leaf mixture
\eqref{eq:tree} within $\sup_Q|\hy-\tilde g|\le C_{d,\beta}\,L\,w\,G^{-1/d}$.
\end{proposition}

The implementation's other prediction rules mix over the root as well as the leaves.
Two of them (\textsc{lik}, \textsc{path}) are universal by the same construction, because
the root's weight is of order $1/G$; the other two (\textsc{prior}, \textsc{posterior})
keep a fixed share on the root of a flat tree, so they are not universal by this
construction (\cref{thm:single,prop:prior}); on deeper trees that share shrinks
(\cref{rem:hierarchy}).

\subsection{Forests}

\begin{lemma}[Stability of composition; \cref{lem:dag}]\label{lem:stable}
Let a finite DAG compute $F$ on a compact $K\subset\R^n$, each vertex $v$ applying a map
$g_v$ that is continuous on an open set containing a closed neighbourhood of its input set
$B_v$ (the values of its inputs as $x$ ranges over $K$).  For every $\varepsilon>0$ there are
$\rho_v,\eta_v>0$ such that replacing every $g_v$ by any map $\hat g_v$ with
$|\hat g_v-g_v|\le\eta_v$ on the closed $\rho_v$-neighbourhood of $B_v$ changes the output by
at most $\varepsilon$ on $K$.
\end{lemma}

\begin{theorem}[Forests of two-input trees are universal approximators]\label{thm:univ-forest}
Let $K\subset\R^n$ be compact, $f:K\to\R$ continuous and $\varepsilon>0$.  There is a
Cobweb forest in which every tree has at most two inputs and passes its activation, whose
output is within $\varepsilon$ of $f$ on $K$.
\end{theorem}

\begin{proof}
By the Stone--Weierstrass theorem \citep{stone1948} the polynomials in $x_1,\dots,x_n$
are dense in $C(K)$: they form an algebra of continuous functions that contains the
constants and separates points.  Take a polynomial $p$ with $\max_K|p-f|<\varepsilon/2$.
Every polynomial is computed by a finite DAG whose sources are the coordinates $x_i$ and
constants and whose other vertices apply $(a,b)\mapsto a+b$, $(a,b)\mapsto ab$ or, where
both arguments are the same value, $a\mapsto a^2$ (by induction on the expression; if $p$ is
a coordinate or a constant, add one identity vertex).  By \cref{lem:stable} with tolerance $\varepsilon/2$ there
are $\rho_v,\eta_v$ for this DAG.  Each vertex map is continuous on the compact set
$K_v^{\rho_v}$, the closed $\rho_v$-neighbourhood of its input set, so by
\cref{thm:univ-tree} a Cobweb tree with two inputs (one, for $a\mapsto a^2$) approximates it within $\eta_v$
there;
a constant is a tree with no input.  The forest of these trees is within $\varepsilon/2$ of
$p$, hence within $\varepsilon$ of $f$, on $K$ (apply \cref{lem:stable} with a tolerance
slightly below $\varepsilon/2$ to make the inequality strict).
\end{proof}

\begin{proposition}[One input per tree is not enough; \cref{prop:fanin}]\label{prop:fanin-main}
If every tree of a forest has at most one input, the output is a function of a single
coordinate.
\end{proposition}

\begin{remark}[An explicit size: the Kolmogorov--Arnold forest; \cref{thm:cka}]\label{rem:ka-main}
The Kolmogorov--Arnold representation
\begin{equation}\label{eq:ka}
  f(x)=\sum_{q=0}^{2n}\Phi_q\Bigl(\sum_{p=1}^n\varphi_{q,p}(x_p)\Bigr)
\end{equation}
is a DAG of
$n(2n+1)$ inner maps, $(2n+1)(n-1)$ additions, $2n+1$ outer maps and $2n$ additions.  The
same argument turns it into a forest of $4n^2+4n$ trees with at most two inputs each, a
size that does not depend on $f$ or $\varepsilon$ (only the trees' leaves do).
\end{remark}

Two further facts place these results.  Trees that pass \emph{labels} are also universal:
trees that categorize each coordinate, feeding a lookup tree, approximate any $f$, but the
lookup tree needs one leaf per cell of the product partition (\cref{thm:label}).  And
merging trees under one common root does not compose them: the merged tree predicts a
convex combination of their predictions (\cref{prop:interweave}).  Composition is what a
forest adds to a tree.

\subsection{What composition buys}

Universality says nothing about size.  Many functions of interest are compositions of
functions of few variables, and for them a forest escapes the dimension.  Call $f$ a
\emph{$k$-composition of depth $D$} if it is computed by a DAG whose sources are the
coordinates and whose every other vertex $v$ applies a map $g_v$ of $k_v\le k$ of its
predecessors.  The \emph{depth} $d(v)$ of a vertex is the number of non-source vertices on the
longest path from a source to $v$, and $D=\max_vd(v)$.  Neighbourhoods are in the max-norm:
$B^{r}=\{u:\operatorname{dist}_\infty(u,B)\le r\}$.

\begin{theorem}[Composition beats dimension]\label{thm:compose}
Let $f$ be a $k$-composition of depth $D$ on a compact $K\subset\R^n$, let $B_v$ be the
set of input vectors of $g_v$ as $x$ ranges over $K$, and suppose that, for some $r>0$,
every $g_v$ is $L$-Lipschitz in the max-norm on $B_v^r$, which lies in a cube $Q_v$ of side
$w$.  Give every vertex the tree of \cref{prop:rate} for the McShane extension of $g_v$ to
$Q_v$, with $m^{k_v}\le G=m^k$ leaves, and let $\eta=C\,L\,w/m$ ($C=\max_{j\le k}C_{j,\beta}$)
be the resulting bound.  If the error of every vertex stays at most $r$---guaranteed when
$\eta\sum_{l<D}L^l\le r$, and checkable after the fact otherwise---then every vertex has
error at most $\eta\sum_{l<d(v)}L^l$, so
\[
  \sup_K|\hat F-f|\;\le\;C\,L\,w\Bigl(\sum_{l=0}^{D-1}L^l\Bigr)G^{-1/k}.
\]
The exponent is set by the fan-in $k$, not by $n$; the constant grows with the depth, which
is at least $\log_kn$ when $f$ depends on all $n$ coordinates.  The corresponding bound for
one tree over all $n$ coordinates is $C_{n,\beta}\,L_f\,w\,G^{-1/n}$ (\cref{prop:rate}).
\end{theorem}

\begin{proof}
Write $e_v=\sup_K|\hat u_v-u_v|$ for the error of the value at vertex $v$ ($e=0$ at the
sources), and $\tilde g_v$ for the extension.  Suppose every predecessor of $v$ has error at
most $e\le r$.  Then $v$'s
approximate input vector is within $e$ (in max-norm) of its exact input vector, so it lies
in $B_v^r\subseteq Q_v$, where the tree is within $\eta$ of $\tilde g_v$ and $\tilde g_v=g_v$
is $L$-Lipschitz:
\[
  |\hat g_v(\hat u)-g_v(u)|\le|\hat g_v(\hat u)-\tilde g_v(\hat u)|+|\tilde g_v(\hat u)-\tilde g_v(u)|\le\eta+Le.
\]
By induction on $d(v)$, $e_v\le\eta\sum_{j<d(v)}L^j$, which is at most $r$ by hypothesis.  At the
output, $d(v)=D$.
\end{proof}

The bound for one tree is only an upper bound, so the theorem does not prove that a single
tree \emph{must} do worse; the measurements of \cref{sec:measure} show that it does.  The same dichotomy between compositional and generic targets is
known for deep and shallow networks \citep{poggio2017}.

%=====================================================================
\section{Learning}\label{sec:learn-main}
%=====================================================================

The constructions above place leaves by hand.  We now let Cobweb learn.

\subsection{One tree}

\begin{theorem}[Cobweb's learning is consistent; \cref{lem:leaves,thm:consistency}]\label{thm:learn-tree}
None of Cobweb's operators deletes a leaf, and with exact matching and no depth limit
every example ends in a leaf that holds only copies of it.  Hence, whatever
hierarchy category utility builds and in whatever order the examples arrive, the
count-weighted leaf mixture \eqref{eq:tree} of the learned tree \emph{equals} the
Nadaraya--Watson kernel regression on all examples seen.  Consequently, for i.i.d.\ examples with bounded $Y$:
if $h\to0$ and $Nh^d\to\infty$, the learned tree converges to $\E[Y\mid X]$ in $L_1$, almost
surely \citep{devroye1989}; and if, on a neighbourhood of a compact set $I$ interior to the
support, $\E[Y\mid X=x]$ is twice differentiable with bounded second derivatives and $X$ has a
continuously differentiable density bounded below, then with $h\propto N^{-1/(d+4)}$ its
root-mean-squared error on $I$ is $O(N^{-2/(d+4)})$: the bias is $O(h^2)$ and the variance
$O(1/(Nh^d))$ uniformly on $I$ \citep{wand1995,gyorfi2002}.
\end{theorem}

The hierarchy plays no role in consistency; it matters for search under a budget, for the
default \textsc{path} rule, and for the labels a tree passes on (\cref{rem:structure}).  The
same fact---nothing stored is deleted---means a tree does not forget: learning new examples
changes old predictions only through the kernel tails of the new leaves.

\subsection{Forests}

A tree in a forest learns to predict its \emph{outcome} $V_\tau$ from its inputs $U_\tau$.
By \cref{thm:learn-tree} it learns $\E[V_\tau\mid U_\tau]$, so a forest has a natural limit,
its \emph{ideal forest}: in topological order every tree emits $m_\tau=\E[V_\tau\mid U_\tau]$,
with $U_\tau$ built from its parents' ideal messages; the ideal output is $F$.  The
\emph{deficit} $\Delta=\E(f-F)^2$, with $f=\E[Y\mid X]$, measures what the assignment of
outcomes loses (\cref{def:ideal}).

\begin{theorem}[A learned forest converges to its ideal forest; \cref{thm:ideal,cor:iff}]\label{thm:learn-forest}
Consider a forest whose trees pass activations, trained level by level on separate parts
of the data (the split protocol).  Under regularity conditions (\cref{ass:regular}) its
output converges uniformly, in probability, to $F$, and its mean squared error to
$\Delta+\E\Var(Y\mid X)$.  It is consistent if and only if $\Delta=0$.
\end{theorem}

\begin{theorem}[Learned forests converge at the rate of their fan-in]\label{thm:learn-rate}
Consider a forest trained with the split protocol, with $D$ levels below the output, and
let $Q$ be the support of $X$.  Write $U_\tau$ and $\hat U_\tau$ for the ideal and learned
inputs of tree $\tau$ (both functions of the raw coordinates $X_{R(\tau)}$ upstream of
$\tau$), $m_\tau=\psi_\tau(U_\tau)$ for its ideal message and
$e_\tau=\sup_{x\in Q}|\hat m_\tau(x)-m_\tau(x)|$ for its error.  Assume:
\begin{enumerate}[label=(\roman*),leftmargin=*,noitemsep]
\item (sufficiency) $\E[V_\tau\mid X_{R(\tau)}]=\psi_\tau(U_\tau)$ for every tree;
\item (smoothness) every $\psi_\tau$ is $L$-Lipschitz in the max-norm on a set containing
  all ideal and learned inputs;
\item (single-tree rate) given its frozen parents, every tree is within $r_N$ of
  $\phi_\tau(\hat u)=\E[V_\tau\mid\hat U_\tau=\hat u]$ uniformly on $\hat U_\tau(Q)$, with
  probability tending to one.
\end{enumerate}
Then, with probability tending to one,
\[
  e_{\text{out}}\;\le\;r_N\sum_{l=0}^{D}(2L)^l .
\]
In particular the output inherits the rate $r_N$ of one tree with $k$ inputs---$N^{-2/(k+4)}$
under the conditions of \cref{thm:learn-tree}---whatever the number $n$ of raw inputs.
Funnel forests on the targets of \cref{prop:tasks}(a) satisfy (i), and (ii) when the
monotone maps of \cref{prop:suffclass} have derivatives bounded away from zero.
\end{theorem}

\begin{proof}
Fix a tree $\tau$ and suppose every parent has error at most $e$, so $\|\hat U_\tau(x)-
U_\tau(x)\|_\infty\le e$ for every $x\in Q$.  Since $\hat U_\tau$ is a function of $X_{R(\tau)}$
(the parents are frozen), the tower property and (i) give
$\phi_\tau(\hat u)=\E[\E[V_\tau\mid X_{R(\tau)}]\mid\hat U_\tau=\hat u]=\E[\psi_\tau(U_\tau)\mid\hat
U_\tau=\hat u]$, so by (ii) $|\phi_\tau(\hat u)-\psi_\tau(\hat u)|\le Le$.  For $x\in Q$, with
$\hat u=\hat U_\tau(x)$ and $u=U_\tau(x)$,
\[
  |\hat m_\tau(x)-m_\tau(x)|\le|\hat\psi_\tau(\hat u)-\phi_\tau(\hat u)|
  +|\phi_\tau(\hat u)-\psi_\tau(\hat u)|+|\psi_\tau(\hat u)-\psi_\tau(u)|\le r_N+2Le
\]
on the event of (iii).  Trees at the lowest level read only raw coordinates ($e=0$), so by
induction a tree at level $l$ has error at most $r_N\sum_{j\le l}(2L)^j$; the finitely many
events of (iii) hold together with probability tending to one.  For (i) on the targets of
\cref{prop:tasks}(a): the proof of \cref{prop:suffclass} shows
$\E[Y\mid X_{R(\tau)}]=H_S(s_S)$, a function of the tree's ideal inputs.
\end{proof}

Hypothesis (iii) is the single-tree rate of \cref{thm:learn-tree} applied to the learned
inputs, whose law must satisfy that theorem's regularity conditions; we assume it rather than
derive it.  The factor $(2L)^l$ is the price of depth: errors compound through levels, which
is why forests should be shallow (\cref{sec:recipe}).

A KAN or a neural network learns inner functions from the error of their composition; a
Cobweb tree receives no error and can only learn to predict something it will observe.  The
natural observable outcomes are the target $Y$ (\emph{funnel}), a part of the input the tree
does not see (\emph{cross-view}), or none, in which case the tree passes labels
(\emph{categorization}).

\begin{proposition}[Which outcomes recover the target; \cref{sec:tasks}]\label{prop:tasks}
\leavevmode
\begin{enumerate}[label=(\alph*),leftmargin=*,noitemsep]
\item If $X$ has compact support and $Y=\Phi\bigl(\sum_jg_j(X_{B_j})\bigr)+\xi$ with $\Phi$
  continuous and strictly increasing, continuous $g_j$, independent blocks $X_{B_j}$ and
  centred integrable noise $\xi$ independent of $X$, every
  funnel forest whose graph is a rooted tree reading each block in exactly one tree has
  $\Delta=0$ (\cref{prop:suffclass}).
\item If $Y=X_1X_2$ with independent inputs symmetric about $0$ and of finite variance, the funnel forest over two
  one-input trees has $F\equiv0$ and $\Delta=\Var Y$: it learns nothing (\cref{prop:xy}).
\item For bounded $Y$, label forests trained with the split protocol, whose output tree
  matches exactly and predicts with the count-weighted leaf mixture, are consistent for every
  target whenever their cells shrink, stay few compared with the data, and their smoothing
  vanishes fast enough (\cref{thm:labelcons}).
\item In a Gaussian two-view model, cross-view forests have $\Delta=0$ when the target
  depends only on what the views share (\cref{prop:multiview}).  Otherwise $\Delta$ can be
  positive; in the Gaussian model it is a closed-form projection, computed for one instance in
  \cref{sec:measure}.
\end{enumerate}
\end{proposition}

\begin{proposition}[No tree can tell its role from its own data;
\cref{thm:locality,cor:parity}]\label{prop:local}
Let $X_1,X_2,X_3$ be independent and uniform on $[-1,1]$, and observe $(X_1,X_2,Y)$ with
either $Y=\operatorname{sign}(X_1X_2)$ or $Y=\operatorname{sign}(X_1X_3)$.  In the first world
$X_2$ is essential, in the second irrelevant, yet $(X_2,Y)$ has the same distribution in
both; any rule that decides, from the data of a tree whose only input is $X_2$, whether to
keep that tree errs with probability at least $1/2$ in one world.
\end{proposition}

What a tree should predict must therefore be decided where trees are composed: choosing
among candidate forests by the validation error of the output is consistent
(\cref{thm:selection}).

%=====================================================================
\section{A recipe for building a forest}\label{sec:recipe}
%=====================================================================

The theory says what a good forest looks like; turning it into a procedure needs only
quantities that Cobweb trees can estimate.  The recipe below plays the role of the KAN
paper's simplification procedure \citep{liu2024kan}.  Each step is justified by a result
above.

\begin{procedure}[Building a Cobweb forest]\label{proc:recipe}
\leavevmode
\begin{enumerate}[leftmargin=*]
\item \emph{Start small.}  One tree per input coordinate (or per given group), composed in
  a balanced binary graph of two-input trees (\cref{thm:univ-forest}: two inputs suffice).
\item \emph{Merge what interacts.}  For every pair of inputs, compare a tree over the pair
  with two one-input trees composed by activations; merge the pair into one tree when the
  merged tree is better by more than a threshold (the \emph{synergy deficit},
  \cref{def:operators}).  Activations lose exactly these interactions
  (\cref{prop:tasks}(b)), and merging keeps the fan-in---and so the rate of
  \cref{thm:compose}---small only when few pairs interact.
\item \emph{Give every tree an observable outcome.}  The target where the parts combine
  monotonically (\cref{prop:tasks}(a)); labels or a merged tree where they interact; another
  view where the target lives in what views share (\cref{prop:tasks}(d)).
\item \emph{Decide at the composing tree.}  Compare candidate forests---always including
  the single merged tree---by the output's validation error (\cref{prop:local},
  \cref{thm:selection}).
\item \emph{Sharpen with data.}  Shrink each tree's variance floor as
  $h\propto N^{-1/(k+4)}$ with $k$ its fan-in (\cref{thm:learn-tree}); this is the analogue of
  a KAN's grid extension.
\item \emph{Prune and read off.}  Drop a tree whose removal does not raise the validation
  error, and plot every remaining tree's activation against its inputs: on targets of
  \cref{prop:tasks}(a) these are the inner functions up to monotone maps.
\end{enumerate}
\end{procedure}

The recipe is greedy and its threshold is a choice.  Pairwise deficits miss interactions of
higher order---for the parity of three signs every pair is independent of the target
(\cref{cor:parity})---which is why step~4 always keeps the merged tree.  On the Feynman formulas and the four-dimensional block target
(\cref{tab:induce}) steps 1--4 recover the groups of interacting variables---$\{B,v,\theta\}$ in
$q(E_f+Bv\sin\theta)$, the two angles of $I.29.16$, the two blocks of the 4-D target---and the
selected forest is within one percent of the best candidate in every case but one; where no
pair interacts, the recipe returns the single tree.

%=====================================================================
\section{Theory meets measurement}\label{sec:measure}
%=====================================================================

Every result above predicts a number.  \Cref{tab:equiv} lists each prediction next to its
measurement on real Cobweb trees---constructed exactly where the result is about
existence, learned with Cobweb's own procedure where it is about learning.  Learned trees
predict with the count-weighted leaf mixture under full expansion and forests use the split
protocol, the setting of the theorems; the variance floor follows the schedule of
\cref{thm:learn-tree} with its constant chosen on validation data.  Rates are least-squares
slopes in log--log coordinates over the sizes shown in the figures; errors are averaged over
three seeds.  \Cref{app:measure} gives the details.

\IfFileExists{tables/equivariance.tex}{%
\begin{table}[t]
\centering\small
\input{tables/equivariance}
\caption{Each prediction of the theory and its measurement.  ``rate'': exponent of the
error against the number of leaves $G$ or of examples $N$.  ``floor'': limit of the
error divided by $\mathrm{sd}(f)$.  ``identity'': largest deviation from an identity that
holds exactly.}
\label{tab:equiv}
\end{table}}{}

\IfFileExists{figs/compose.pdf}{%
\begin{figure}[tp]
\centering
\includegraphics[width=\linewidth]{figs/compose.pdf}
\caption{Composition beats dimension (\cref{thm:compose}).  Sup error of constructed
Cobweb trees and forests on a composition of two-input maps in $n=2,4,8$ dimensions,
against the total number of leaves.  Dashed: the predicted slopes, $-1/2$ for the forest
and $-1/n$ for one tree.}
\label{fig:compose}
\end{figure}}{}

\IfFileExists{figs/learn.pdf}{%
\begin{figure}[t]
\centering
\includegraphics[width=\linewidth]{figs/learn.pdf}
\caption{Learning.  Left: one learned tree, interior RMSE against $N$ in $d=1,2,3$
dimensions; dashed: the predicted slopes $-2/(d+4)$.  Right: learned funnel forest and one
learned tree on $\Phi(\sum_j g(X_{B_j}))$ with two-dimensional blocks, $n=2,4,8$.}
\label{fig:learn}
\end{figure}}{}

\paragraph{Representation.}
One constructed tree's sup error falls with the predicted exponent $-1/d$
(measured $-1.06$, $-0.51$, $-0.30$ for $d=1,2,3$).  On the composition of two-input maps
(\cref{fig:compose}) the forest's exponent does not degrade with the dimension ($-0.60$,
$-0.47$, $-0.55$ for $n=2,4,8$, against the predicted $-0.5$), while the single tree's exponent shrinks with $n$ as
predicted ($-0.54$ and $-0.23$ against $-0.5$ and $-0.25$; at $n=8$ its grids, with at most five
points per axis, are still pre-asymptotic: $-0.06$ against $-0.125$).  The difference in size
is large: in eight dimensions a forest with $28{,}672$ leaves in all has sup error $0.020$, a
single tree with $390{,}625$ leaves $0.22$.  With smooth maps the forest converges faster than the Lipschitz bound ($-0.74$ to
$-0.92$), as an upper bound allows; the single tree at $n=8$ is again pre-asymptotic.
The a priori margin condition of \cref{thm:compose} would need far finer trees, but its a
posteriori form can be checked: in every forest with at least $64$ leaves per tree, every
approximate input stayed inside its margin at every level, so these runs satisfy the
theorem's hypothesis, and the forest exponents are fitted on them.

\paragraph{Learning.}
A learned tree's count-weighted leaf mixture equals the Nadaraya--Watson
estimate to within $3\times10^{-7}$, float precision, and its interior error falls at $-0.41$,
$-0.33$, $-0.31$ for $d=1,2,3$ against the predicted $-0.40$, $-0.33$, $-0.29$.  For learned
forests \cref{thm:learn-rate} predicts the rate of one tree with $k=2$ inputs,
$N^{-2/(k+4)}=N^{-1/3}$, in every dimension; the forests follow it ($-0.36$ for $n=4$, $-0.38$
for $n=8$).  One tree
over all inputs converges at least as fast as its exponent $-2/(n+4)$ for $n=4$ ($-0.32$) and
is still pre-asymptotic for $n=8$ ($-0.12$).  In eight dimensions at $N=16{,}000$ the forest's
normalised error is $0.25$ and the single tree's $0.74$ (\cref{fig:learn}).

\paragraph{What trees predict.}
Each kind of outcome behaves as \cref{prop:tasks} says.  On $x_1x_2$ the funnel forest's
normalised error is $1.00$ at every $N$, the predicted floor: its trees' activations are
constant.  The label forest, with a small smoothing $\alpha=10^{-4}$ in its output tree, falls from
$0.24$ to $0.13$ to $0.083$ as $N$ grows from $10^3$ to $1.6\times10^4$, and the validated label
depth grows with $N$ ($3$, $3$, $4$): its cells refine, the behaviour \cref{thm:labelcons}
requires.  These are finite-$N$ illustrations of the theorem, whose conditions are
asymptotic.  With the default smoothing $\alpha=0.1$, for which $\alpha M_N$ does not vanish,
the forest stalls near $0.6$.  On the
view-specific Gaussian target the cross-view forest approaches its predicted floor
$\sqrt{\Delta/\Var f}=0.696$ from above ($0.81$, $0.78$, $0.72$).

%=====================================================================
\section{Case studies}\label{sec:cases}
%=====================================================================

The theorems are constructive.  This section runs the constructions on concrete targets with
real Cobweb trees, so that every step of a proof becomes a measurement.

\subsection{The universality proof, executed}

\Cref{thm:univ-forest} approximates $f$ by a polynomial and then replaces each operation of
the polynomial by a Cobweb tree.  We run it with the Bernstein polynomials
$B_mf(x,y)=\sum_{i,j\le m}f(\tfrac im,\tfrac jm)\,b_{m,i}(x)\,b_{m,j}(y)$, the classical constructive
proof of the Weierstrass theorem, for $f(x,y)=\exp(\sin\pi x+y^2)$ on $[0,1]^2$.  One-input trees
compute the basis functions $b_{m,i}$, two-input trees compute the products
$f(\tfrac im,\tfrac jm)\,ab$ and a balanced binary tree of two-input trees adds them; every tree is
built as in \cref{thm:univ-tree}, from $G$ leaves on the box of its inputs.  The proof makes
two predictions: the forest converges to $B_mf$ as its trees refine (\cref{lem:stable}), with
a tolerance that tightens as the forest grows; and $B_mf$ converges to $f$, at the rate $1/m$
for smooth $f$.  \Cref{fig:bernstein} shows both.  With $64^2$ leaves per tree the forest is
within $0.010$ of $B_2f$ (23 trees) and within $0.050$ of $B_{12}f$ (363 trees), and larger
forests need finer trees, as the lemma says.  Its error to $f$ follows the Bernstein error
down, from $2.22$ at $m=2$ to $0.59$ at $m=12$ (local slope $-0.73$, approaching $-1$).  The
forest never does better than the polynomial it implements: the proof's two errors add, and
the measurements separate them.

\IfFileExists{figs/bernstein.pdf}{%
\begin{figure}[t]
\centering
\includegraphics[width=\linewidth]{figs/bernstein.pdf}
\caption{The proof of \cref{thm:univ-forest}, executed with Bernstein polynomials.  Left: sup
distance between the Cobweb forest and the polynomial $B_mf$ it implements, against the leaves
per tree.  Right: the forest's error to $f$ (with $64^2$ leaves per tree) and the Bernstein
error $\sup|B_mf-f|$, against the degree $m$; dashed, slope $-1$.}
\label{fig:bernstein}
\end{figure}}{}

\subsection{Physics formulas}

\Cref{thm:compose} says that a target which is a composition of few-input maps should be
represented by a forest that mirrors the composition.  Physics formulas are such targets.  We
compile nine formulas of the Feynman benchmark \citep{liu2024kan}, with two to nine inputs, into
forests: every operation of the formula's expression graph becomes one constructed Cobweb tree
with at most two inputs, built on the box of its inputs with the Lipschitz extension of
\cref{thm:compose}.  We compare each forest with one constructed tree over all inputs
(\cref{tab:feynman-cases}).

Three observations match the theory.  First, every forest converges as its trees refine, as
\cref{thm:compose} guarantees.  Second, where the formula has more inputs than any one of its operations reads,
the forest is more accurate at a comparable number of leaves: $0.0035$ against $0.038$ on
$I.50.26$ ($n=4$), $0.039$ against $0.32$ on $II.6.15a$ ($n=6$), and $0.0030$ against
$0.0063$ on $I.6.20b$.  On the nine-input formula $I.9.18$ the forest has overtaken the single
tree ($0.15$ against $0.17$) and is still falling fast ($1.18$, $0.59$, $0.15$ over its last
three refinements).  Third, the theorem's constant predicts where forests struggle.  $\kappa$ is
computed from the formula alone, before any tree is built.  It has rank correlation $0.80$
with the forest's error across the nine formulas.  The two formulas with by far the largest
constant---$I.29.16$, a square root of a near-cancelling sum, and $I.13.12$, products of large
values with differences of reciprocals---are the two on which the forest loses.  The forest
also loses on $I.6.20$, which has only two inputs, so there is no dimension to beat.  The
lesson for building forests (\cref{sec:recipe}) is to choose graphs whose operations have small
Lipschitz constants.

\IfFileExists{tables/feynman_cases.tex}{%
\begin{table}[t]
\centering\small
\resizebox{\linewidth}{!}{\input{tables/feynman_cases}}
\caption{Feynman formulas compiled into Cobweb forests (one constructed tree per operation of
the expression graph, fan-in at most two) and one constructed tree over all $n$ inputs.
Errors are RMSE$/\mathrm{sd}(f)$ on 500 uniform points of the formula's input box.  $\kappa$:
the constant of \cref{thm:compose} for this graph, $\sum_vL_vw_vA_v/\mathrm{sd}(f)$, with $L_v$
the Lipschitz constant of operation $v$, $w_v$ the width of its input box and $A_v$ the
largest product of Lipschitz constants from $v$ to the output.}
\label{tab:feynman-cases}
\end{table}}{}

\subsection{No forgetting, and reading a forest}

Two consequences of the theory are visible directly.  Because no stored example is deleted
(\cref{thm:learn-tree}), a tree trained on the five-peak task of \citet[Fig.~3.4]{liu2024kan}, one
region at a time, keeps its error on everything seen so far between $0.037$ and $0.048$ through
all five phases (\cref{fig:continual}).  And because every tree of a funnel forest learns a
conditional expectation, its trees can be plotted and read: on $\exp(\sin\pi x_1+x_2^2)$ the
two input trees learn $e^{\sin\pi x_1}$ and $e^{x_2^2}$ up to constants (\cref{fig:interpret}),
the inner functions of the target up to the monotone maps of \cref{prop:suffclass}.

%=====================================================================
\section{Related work}\label{sec:related}
%=====================================================================

\paragraph{Universal approximation.}  One-hidden-layer networks approximate every
continuous function on a compact set \citep{cybenko1989,hornik1989}; the
Kolmogorov--Arnold theorem \citep{kolmogorov1957,arnold1957,lorentz1966,sprecher1965}
represents every continuous function through one-variable functions and addition, and
KANs \citep{liu2024kan} make that representation learnable.  Deep networks avoid the curse of
dimensionality on compositional targets \citep{poggio2017}.  Our theorems are the
counterparts for a learner without gradients, and our proof of universality follows the
Stone--Weierstrass route \citep{stone1948}.

\paragraph{Kernel smoothing.}  A learned Cobweb tree's leaf mixture is a Nadaraya--Watson
estimate (\cref{thm:learn-tree}), so its consistency and rates come from that literature
\citep{devroye1989,gyorfi2002,wand1995}.  What Cobweb adds is incremental structure---the
hierarchy, its labels, its search---and what the forest adds is composition.

\paragraph{Concept formation.}  Cobweb \citep{fisher1987} and its information-theoretic
category utility \citep{corter1992,maclellan2026rational} learn concept hierarchies
incrementally.  Systems that let one tree read another's concepts
\citep{maclellan2021conv,singaravadivelan2026trellis} motivated the forests studied here; we
use none of their results.

%=====================================================================
\section{Conclusion}\label{sec:conclusion}
%=====================================================================

Cobweb is a universal approximator: one tree suffices to represent any continuous
function, forests of two-input trees suffice as well, and forests represent compositional
functions at a rate set by their fan-in rather than by the dimension.  Cobweb's own
learning is consistent for a tree, and a learned forest converges to an explicit limit
whose distance from the target is decided by what each tree predicts.  Every one of these
statements predicts a number that the implementation reproduces.

\paragraph{Limitations.}
The convergence theorems assume the count-weighted leaf mixture, exact matching and, for
forests, the split protocol; \cref{thm:ideal} assumes a uniform law for kernel smoothers up
to the boundary ((A5) of \cref{ass:regular}).  The single-tree rates are upper bounds, and \cref{thm:learn-rate} assumes the single-tree
rate on the learned inputs (its hypothesis (iii)).  The label conditions
of \cref{thm:labelcons} are not derived from category utility, and the recipe is greedy.

\paragraph{A cognitive note.}
The three outcomes a tree can be given are three things a learner can do with what it
perceives: categorize it, predict an outcome from it, or predict another part of it.  The
results say when each suffices, and \cref{prop:local} is a precise form of the
credit-assignment problem \citep{minsky1961}: a module cannot learn its role from what it
alone sees.

\appendix
\numberwithin{theorem}{section}
